\documentclass[12pt]{article}
\usepackage[T1]{fontenc}
\usepackage[french]{babel}
\usepackage{datetime}
\usepackage{subcaption}
\usepackage[letterpaper, margin=1.9cm]{geometry}
\usepackage{graphicx}
\usepackage{epstopdf}
\usepackage{siunitx}
\sisetup{output-decimal-marker = {,}}
\graphicspath{ {../figures} }
\usepackage{derivative}
\usepackage{booktabs}
\usepackage[flushleft]{threeparttable}
\usepackage{amsmath}

\usepackage{csquotes}
\usepackage[
backend=biber,
style=ieee,
%style=verbose,
citestyle=ieee,
%sorting=nyt,
%url=false
]{biblatex}
\bibliography{references}

\title{MEC6616 --- LAP2}


\author{Aurélie Grebe\\Noémie Herbinière}
\newdate{duedate}{14}{09}{2026}
\date{\displaydate{duedate}}

\begin{document}
\maketitle

\section{Introduction}

Dans cette œuvre, nous présentons un solveur qui résout l'équation
d'advection-diffusion en 1-D par volumes finis telle que présentée par Versteeg
et Malalasekera\,\cite{versteeg_introduction_2007}. Nous expliciterons d'abords
quelques détails de l'implémentation avant de comparé des schémas centrés et en
amont à l'aide de deux exemples.

\section{Méthodologie}

On cherche à résoudre l'équation de convection-diffusion:
\begin{equation}
	\odv{}{x}(\rho u \phi) = \odv{}{x}\left( \Gamma\odv{\phi}{x}\right)
\end{equation}
avec la masse volumique du fluide $\rho$, la vitesse $u$, la diffusivité
$\Gamma$ et le variable scalaire $\phi$.

\section{Résultats}

\subsection{Exemple 5.1}

\subsection{Exemple 5.2}
\begin{figure}[]
	\centering
	\begin{subfigure}[t]{0.45\textwidth}
		\includegraphics[width=\textwidth]{ex5_2i.png}
		\caption{Résultats analytiques et par volumes finis de l'exemple 5.2 i)}
		\label{fig:ex5_2}
	\end{subfigure}
	\hfill
	\begin{subfigure}[t]{0.45\textwidth}
		\includegraphics[width=\textwidth]{epsilon5_2i.png}
		\caption{Erreur de solution selon la finesse de discrétisation
		spatiale.}
		\label{fig:err5_2}
	\end{subfigure}
	\caption{}
\end{figure}

\begin{figure}[]
	\centering
	\begin{subfigure}[t]{0.45\textwidth}
		\includegraphics[width=\textwidth]{ex5_2ii.png}
		\caption{Résultats analytiques et par volumes finis de l'exemple 5.2 ii)}
		\label{fig:ex5_2}
	\end{subfigure}
	\hfill
	\begin{subfigure}[t]{0.45\textwidth}
		\includegraphics[width=\textwidth]{epsilon5_2ii.png}
		\caption{Erreur de solution selon la finesse de discrétisation
		spatiale.}
		\label{fig:err5_2}
	\end{subfigure}
	\caption{}
\end{figure}

\printbibliography[
	heading=bibintoc,
	title={Références}
]
	
\end{document}
